The center is $O=(0,0)$ and $OP=65$, while $OQ=OR=25$. The tangent length is $PQ=\sqrt{65^2-25^2}=60$. If $M=OP\cap QR$, symmetry makes $M$ the midpoint of $QR$. In right triangle $OQP$, the altitude to hypotenuse $OP$ is
\[
QM=\dfrac{OQ\cdot QP}{OP}=\dfrac{25\cdot60}{65}=\dfrac{300}{13}.
\]
Hence $\boxed{QR=\dfrac{600}{13}}$.
\begin{asy}
texpreamble("\usepackage{amsmath,amssymb}");
size(240);
defaultpen(fontsize(9pt));
pair O=(0,0),P=(39,52),e=P/65,n=rotate(90)*e,M=(125./13)*e,Q=M+(300./13)*n,R=M-(300./13)*n;draw(circle(O,25));draw(P--Q--O--R--P);draw(Q--R,blue);draw(O--P,dashed);label("$O$",O,(-1,-1));label("$P$",P,(1,1));label("$Q$",Q,(-1,0));label("$R$",R,(1,-1));label("$M$",M,(0,-1));
\end{asy}
